\section*{ID9935: Relation: A\_rm total!,\_μ}
\paragraph{Metadata.}
PiID: 2295524988727 | RTEID: RTE:P38579.8113:T0.7358 | UUID: e689b417-3a27-5d74-9eb0-a4b656445750

\paragraph{Canonical Statement.}
ty, momentum flux, spatial stress) measured by an observer, and instructions to compute/diagonalize principal pressures. 1. Assumptions and notation • Signature: \$(-,+,+,+)\$. Coordinate indices \$\textbackslash{}mu,\textbackslash{}nu=0,1,2,3\$ where \$x\textasciicircum{}0=t\$. Spatial indices \$i,j,k=1,2,3\$. • Metric: \$g\_\{\textbackslash{}mu\textbackslash{}nu\}(t,\textbackslash{}mathbf x)\$ — use the rotating-frame metric \$g\_\{\textbackslash{}mu\textbackslash{}nu\}[\textbackslash{}omega]\$ defined in your notes. Define its inverse \$g\textasciicircum{}\{\textbackslash{}mu\textbackslash{}nu\}\$ and determinant \$g=\textbackslash{}det(g\_\{\textbackslash{}mu\textbackslash{}nu\})\$. • 4-potential: \$A\_\textbackslash{}mu(t,\textbackslash{}mathbf x)=\{\textbackslash{}Phi,,A\_i\}\$ where \$A\_i\$ includes the dipole + any gyroscopic contributions from your file. Label this combined potential as \$A\_\{\textbackslash{}rm total\}!,\{\}\_\textbackslash{}mu\$. • Field tensor: \$F\_\{\textbackslash{}mu\textbackslash{}nu\}=\textbackslash{}partial\_\textbackslash{}mu A\_\textbackslash{}nu-\textbackslash{}partial\_\textbackslash{}nu A\_\textbackslash{}mu\$. • Units: SI-like with factors \$\textbackslash{}varepsilon\_0,\textbackslash{}mu\_0,c\$ shown explicitly when helpful; geometrized units can be obtained by setting \$\textbackslash{}varepsilon\_0=\textbackslash{}mu\_0=1\$ and \$c=1\$. 2. Covariant electromagnetic stress–energy tensor Start from the standard expression T\textasciicircum{}\{\textbackslash{}mu\textbackslash{}nu\}\_\{\textbackslash{}rm EM\} \textbackslash{},=\textbackslash{}, F\textasciicircum{}\{\textbackslash{}mu\textbackslash{}alpha\}\{F\textasciicircum{}\{\textbackslash{}nu\}\}\_\{\textbackslash{}!\textbackslash{}alpha\} - \textbackslash{}tfrac\{1\}\{4\}\textbackslash{},g\textasciicircum{}\{\textbackslash{}mu\textbackslash{}nu\}F\textasciicircum{}\{\textbackslash{}alpha\textbackslash{}beta\}F\_\{\textbackslash{}alpha\textbackslash{}beta\}. Steps to form components symbolically: • Compute \$F\_\{\textbackslash{}mu\textbackslash{}nu\}\$

\paragraph{Canonical Equation.}
\[
A_{\rm total}!,{}_\mu
\]

\paragraph{Dictionary.}
Relation: A\_rm total!,\_μ — A\_rm total!,\_μ

\paragraph{Encyclopedia.}
Relation: A\_rm total!,\_μ is a symbolic expression, written A\_rm total!,\_μ. It introduces a symbol or relation used elsewhere in the framework. It is built from μ, A\_ total. Within the Trawin Zero Point registry it serves as one element of the framework's mathematical narrative.
